\section{A completed complex-power sine transform}
\label{sec:beta}

Put
\begin{equation}\label{beta:weight}
 \ell(x)=\log\frac{\sin\pi x}{\pi},\qquad L=\log(2\pi),\qquad
 W_\lambda(x)=e^{2\lambda\ell(x)}.
\end{equation}
The normalization by $\pi$ makes the leading endpoint coefficient equal
to one. It is useful both for coordinate finite parts and for keeping
the later crossing subtraction unambiguous. Define
\begin{equation}\label{beta:A}
 A(s,\lambda)=\int_0^1\zeta(s,x)W_\lambda(x)\dd x,
 \qquad
 C_0(\lambda)=\frac{\Gamma(1+2\lambda)}
 {\Gamma(1+\lambda)^2(2\pi)^{2\lambda}}.
\end{equation}
An initial domain sufficient for everything in the first theorem is
$\Re\lambda>-1/2$, $\Re s<0$. All subsequent values of $A$ mean its
meromorphic continuation.

\subsection{Fourier coefficients and a polylogarithmic Euler integral}

\begin{theorem}[The complex-power transform]\label{beta:transform}
For $\Re\lambda>-1/2$ the Fourier coefficients of $W_\lambda$ are
\begin{equation}\label{beta:cn}
 c_n(\lambda)=\int_0^1W_\lambda(x)e^{-2\pi inx}\dd x
 =\frac{(-1)^n\Gamma(1+2\lambda)}
 {(2\pi)^{2\lambda}\Gamma(1+\lambda+n)\Gamma(1+\lambda-n)}.
\end{equation}
They satisfy $c_{-n}=c_n$, $c_0=C_0$, and, for $n\geq1$,
\begin{equation}\label{beta:pochhammer}
 c_n(\lambda)=C_0(\lambda)\frac{(-\lambda)_n}{(1+\lambda)_n}.
\end{equation}
Consequently,
\begin{equation}\label{beta:dirichlet}
 A(s,\lambda)=B(s)\sum_{n=1}^{\infty}c_n(\lambda)n^{s-1},\qquad
 B(s)=2\Gamma(1-s)(2\pi)^{s-1}\sin\frac{\pi s}{2}.
\end{equation}
In the strip $-1/2<\Re\lambda<0$, $\Re s<0$, a second representation is
\begin{align}\label{beta:euler}
 A(s,\lambda)
 &=-\frac{2}{\pi}\Gamma(1-s)(2\pi)^{s-1-2\lambda}
      \sin\frac{\pi s}{2}\sin(\pi\lambda)\notag\\
 &\hspace{2em}\times
 \int_0^1t^{-\lambda-1}(1-t)^{2\lambda}\Li_{1-s}(t)\dd t.
\end{align}
Both formulas continue meromorphically wherever their two sides are
defined by continuation.
\end{theorem}

\begin{proof}
Reflection about $1/2$ kills the sine coefficient. If
$I_n=\int_0^\pi\sin^{2\lambda}t\cos(2nt)\dd t$, integration of
$\frac{\mathrm d}{\mathrm dt}
 [\sin^{2\lambda+1}t\cos((2n+1)t)]$ gives
\[
 (\lambda-n)I_n+(\lambda+n+1)I_{n+1}=0.
\]
The boundary term vanishes for $\Re\lambda>-1/2$.
The beta integral gives
$I_0=\sqrt\pi\,\Gamma(\lambda+1/2)/\Gamma(\lambda+1)$.
The recurrence and duplication formula prove
\eqref{beta:cn}--\eqref{beta:pochhammer}. The apparent exceptions at
integer parameters are read by continuity.

The Hurwitz Fourier expansion for $\Re s<0$ is absolutely convergent.
Its sine terms integrate to zero and its cosine terms give
\eqref{beta:dirichlet}; see also \cite[\S25.11]{DLMFHurwitz}.
The decay $c_n=O(n^{-1-2\Re\lambda})$ shows that the resulting
Dirichlet series itself converges absolutely for
$\Re s<1+2\Re\lambda$, away from exceptional parameters interpreted
by continuity. Lastly, Gamma reflection gives
\[
 c_n(\lambda)=-\frac{\sin\pi\lambda}{\pi(2\pi)^{2\lambda}}
 \int_0^1t^{n-\lambda-1}(1-t)^{2\lambda}\dd t.
\]
Absolute convergence in the stated strip permits summation under the
integral, producing \eqref{beta:euler}.
\end{proof}

These Fourier and beta mechanisms are classical. In particular,
Espinosa and Moll treat Hurwitz transforms of trigonometric powers
and the first log-sine transform \cite[\S\S5,13]{EspinosaMoll}.
The purpose of Theorem~\ref{beta:transform} here is to supply a common
generator whose \emph{completed germ at the joint resonance} can be
specified at every Stieltjes and argument-derivative order.

For an integer $m\geq1$ the Fourier expansion terminates:
\begin{equation}\label{beta:integer}
 A(s,m)=\frac{B(s)}{(2\pi)^{2m}}
 \sum_{n=1}^m(-1)^n\binom{2m}{m+n}n^{s-1}.
\end{equation}
For example, comparison of the constant term at $s=1$ gives the
ordinary convergent identity
\begin{equation}\label{beta:psiinteger}
 \int_0^1\psi(x)W_m(x)\dd x
 =-C_0(m)(\gamma+L)
 +\frac{2}{(2\pi)^{2m}}
   \sum_{n=1}^m(-1)^n\binom{2m}{m+n}\log n.
\end{equation}
Higher Laurent coefficients give the analogous all-index Stieltjes
moments as finite Gamma and logarithmic expressions.

\subsection{The complete crossing polar term}

Let $P_p(u)=(1+u)_p$, with $P_0=1$, and define the entire-in-$\lambda$
Taylor coefficients
\begin{align}\label{beta:d}
 d_j(\lambda)
 &=[x^{2j}]\left(\frac{\sin\pi x}{\pi x}\right)^{2\lambda}\notag\\
 &=[z^j]\exp\left(-2\lambda
                 \sum_{r=1}^{\infty}\frac{\zeta(2r)}r z^r\right).
\end{align}
Thus $d_0=1$, $d_1=-2\lambda\zeta(2)$, and
$d_2=-\lambda\zeta(4)+2\lambda^2\zeta(2)^2$.

\begin{theorem}[All logarithmic Stieltjes moments]\label{beta:completion}
For every integer $p\geq0$ the expression
\begin{align}\label{beta:Hp}
 \mathcal H_p(u,\lambda)
 ={}&(-1)^pP_p(u)A(1+p+u,\lambda)
       -\ind_{p=0}\frac{C_0(\lambda)}u\notag\\
 &-\ind_{2\mid p}
       \frac{(-1)^pP_p(u)d_{p/2}(\lambda)}{2\lambda-u}
\end{align}
extends holomorphically to a neighborhood of $(u,\lambda)=(0,0)$.
For all integers $m,k\geq0$,
\begin{equation}\label{beta:allmoments}
 \boxed{\displaystyle
 \FP\int_0^1\gamma_m^{(p)}(x)\ell(x)^k\dd x
  =\frac{(-1)^m m!k!}{2^k}
      [u^m\lambda^k]\mathcal H_p(u,\lambda).}
\end{equation}
The finite part is exactly the unit-coordinate convention
\eqref{scope:FP}.
\end{theorem}

\begin{proof}
The Hurwitz shift equation gives, near $x=0$,
\begin{equation}\label{beta:local}
 \partial_x^p\left(\zeta(1+u,x)-\frac1u\right)
 =(-1)^pP_p(u)x^{-1-p-u}+h_p(u,x),
\end{equation}
where $h_p$ is holomorphic jointly near $(u,x)=(0,0)$.
For $p=0$ the subtraction $1/u$ is retained inside $h_0$;
for $p>0$ its derivative is zero. Moreover
\[
 W_\lambda(x)=x^{2\lambda}\sum_{j\geq0}d_j(\lambda)x^{2j}.
\]
Integration of the singular product from zero to a fixed
$\delta\in(0,1)$ produces the denominators
$2\lambda-u+2j-p$. Only $j=p/2$ can vanish at the origin, and it
occurs only when $p$ is even. Its full numerator is the one displayed
in \eqref{beta:Hp}. Subtracting finitely many endpoint terms makes
the remaining integral locally normally convergent; all the other
denominators are bounded away from zero. The endpoint $x=1$ has only
the integrable analytic family $(1-x)^{2\lambda}$ times analytic
functions and creates no divisor through $(0,0)$. The pole $1/u$ of
the original Hurwitz function contributes $C_0(\lambda)/u$ only
when $p=0$. This proves joint holomorphy.

For completeness, the local resonant integral is
$N(u,\lambda)\delta^{2\lambda-u}/(2\lambda-u)$. After subtracting
the entire quotient $N(u,\lambda)/(2\lambda-u)$, its Taylor
coefficients are those of
$N(u,\lambda)(\delta^{2\lambda-u}-1)/(2\lambda-u)$.
These are precisely the constants left by integrating the corresponding
powers of $\log x$ in the coordinate $x$. All nonresonant terms have
the same property by ordinary analytic continuation of
$(\delta^a-\epsilon^a)/a$. Thus Taylor coefficients commute with the
specified finite part after the full subtraction. Applying
\eqref{scope:stieltjes} and
$W_\lambda=\sum_k2^k\ell^k\lambda^k/k!$ proves
\eqref{beta:allmoments}.
\end{proof}

\begin{corollary}\label{beta:zero}
For every $m,p\geq0$,
\[
 \FP\int_0^1\gamma_m^{(p)}(x)\dd x=0.
\]
\end{corollary}
\begin{proof}
The meromorphic identity $A(s,0)=0$ follows from
\eqref{beta:dirichlet}, since all nonzero Fourier coefficients
vanish. Also $d_j(0)=0$ for $j>0$. In \eqref{beta:Hp} the two
remaining terms for $p=0$ cancel along $\lambda=0$; for $p>0$ each
term vanishes there. Hence $\mathcal H_p(u,0)=0$ as a holomorphic
germ.
\end{proof}

\begin{remark}[A removable limit need not equal a coordinate finite part]
\label{beta:anomaly}
For an even integer $p\geq2$, meromorphic continuation from a
convergence domain gives
\begin{equation}\label{beta:anomalyformula}
 \lim_{\lambda\to0}
 \left(\int_0^1\psi^{(p)}(x)W_\lambda(x)\dd x\right)_{\mathrm{cont}}
 =2(p-1)!\zeta(p),
 \qquad
 \FP\int_0^1\psi^{(p)}(x)\dd x=0.
\end{equation}
Indeed $d_{p/2}'(0)=-4\zeta(p)/p$, and the unsubtracted crossing
quotient along $u=0$ is
$-p!d_{p/2}(\lambda)/(2\lambda)$ for the polygamma integral.
This gives the first value, while Corollary~\ref{beta:zero} gives the
second. For $p=2$ the discrepancy is $2\zeta(2)$. The correction is
to remove the \emph{full analytic numerator}, including terms that
vanish at the crossing. This is a normalization warning about a
tempting extension, not a claim that the incoming completed-germ
theorems make this error.
\end{remark}

\subsection{Cotangent powers and normalized antiderivatives}

Let $K(x)=\pi\cot\pi x$. The identity
$K^2=W_{\lambda-1}/W_\lambda-\pi^2$ implies
\begin{equation}\label{beta:evencot}
 \int_0^1\zeta(s,x)K(x)^{2h}W_\lambda(x)\dd x
 =\sum_{j=0}^h\binom hj(-\pi^2)^{h-j}A(s,\lambda-j).
\end{equation}
Since $W_\mu'=2\mu K W_\mu$ and
$\partial_x\zeta(s,x)=-s\zeta(s+1,x)$, integration by parts in
an endpoint-convergent domain gives
\begin{equation}\label{beta:oddcot}
 \int_0^1\zeta(s,x)K(x)^{2h+1}W_\lambda(x)\dd x
 =\sum_{j=0}^h\binom hj(-\pi^2)^{h-j}
   \frac{s}{2(\lambda-j)}A(s+1,\lambda-j).
\end{equation}
These are meromorphic identities. A zero denominator on the right
requires continuation and, for a coordinate finite part, the same
local completion as above. The polynomial cotangent closure already
proved in \cite{RepoExact} is therefore embedded in a complex-power
family rather than claimed afresh.

The symmetric Gamma primitive is particularly simple:
\begin{proposition}\label{beta:gammaprimitive}
For $\Re\lambda>-1/2$,
\begin{equation}\label{beta:gamma}
 \int_0^1\log\Gamma(x)W_\lambda(x)\dd x=-\frac14C_0'(\lambda).
\end{equation}
Consequently every $k\geq0$ satisfies the ordinarily convergent identity
\begin{equation}\label{beta:gammalogs}
 \int_0^1\log\Gamma(x)\ell(x)^k\dd x
 =-\frac{C_0^{(k+1)}(0)}{2^{k+2}}.
\end{equation}
\end{proposition}
\begin{proof}
Gamma reflection gives
$\log\Gamma(x)+\log\Gamma(1-x)=-\ell(x)$. Multiply by the symmetric
weight and integrate. The right side is $-C_0'(\lambda)/2$, while
the left side is twice the desired integral. Differentiation is
justified by endpoint domination on compact subsets of the stated
half-plane.
\end{proof}

For arbitrary Stieltjes index, the normalized antiderivative generator
is
\begin{equation}\label{beta:primitivegenerator}
 \mathcal P(u,x)=-\frac{\zeta(u,x)-\zeta(u,1)+x-1}{u}.
\end{equation}
Its apparent singularity at $u=0$ is removable because
$\zeta(0,x)=1/2-x$. The Hurwitz derivative identity proves
\begin{equation}\label{beta:primitivecoeff}
 \partial_x\mathcal P(u,x)=\zeta(1+u,x)-\frac1u,
 \qquad
 \int_1^x\gamma_m(t)\dd t=(-1)^m m![u^m]\mathcal P(u,x).
\end{equation}
In particular $\mathcal P(0,x)=-\log\Gamma(x)$ by Lerch's formula.
The normalization at $x=1$ fixes the integration constant before any
reflection or finite-part operation. Iterating ordinary integration
of this holomorphic generator gives normalized higher primitives.

\subsection{Harmonic jets of the weight}

\begin{proposition}\label{beta:harmonicjets}
For every integer $n\geq1$ the germ of its Fourier coefficient is
\begin{align}\label{beta:jet}
 c_n(\lambda)=-\frac\lambda n\exp\bigg(&-2L\lambda
 +\sum_{j\geq2}\frac{(-1)^j(2^j-2)\zeta(j)}j\lambda^j\notag\\
 &-\sum_{j\geq1}
 \frac{H_{n-1}^{(j)}+(-1)^{j+1}H_n^{(j)}}j\lambda^j\bigg).
\end{align}
Thus every logarithmic sine moment is generated by finite polynomials
in generalized harmonic numbers and ordinary zeta constants.
\end{proposition}
\begin{proof}
Factor \eqref{beta:pochhammer} as
\[
 c_n=-\frac\lambda n C_0(\lambda)
  \frac{\prod_{j=1}^{n-1}(1-\lambda/j)}
       {\prod_{j=1}^{n}(1+\lambda/j)}.
\]
Expanding the logarithm of each finite product and using the Taylor
series of $\log\Gamma(1+\lambda)$ gives \eqref{beta:jet}.
Each fixed coefficient requires only finitely many terms. Initially
the differentiated Dirichlet series converges normally in a smaller
half-plane; its extension is then supplied by
Theorem~\ref{beta:completion}.
\end{proof}

The first two terms illustrate the inclusive harmonic convention:
\begin{equation}\label{beta:firstjets}
 c_n(\lambda)=-\frac\lambda n+\lambda^2
 \left(\frac{2(L+H_n)}n-\frac1{n^2}\right)+O(\lambda^3).
\end{equation}
More explicitly, if $J_k(s)=\int_0^1\zeta(s,x)\ell(x)^k\dd x$,
then in its initial domain
\begin{equation}\label{beta:belltransform}
 J_k(s)=\frac{k!}{2^k}B(s)
 \sum_{n\geq1}[\lambda^k]c_n(\lambda)n^{s-1}.
\end{equation}
This is a finite Bell-polynomial construction at each fixed $k$, not
an assertion that all resulting harmonic constants have already been
reduced to ordinary ones.

\subsection{The first two logarithmic moments and an exact cancellation}

For $k=1$, the classical transform becomes
\begin{align}\label{beta:J1}
 J_1(s)&=-\Gamma(1-s)(2\pi)^{s-1}
                \sin\frac{\pi s}{2}\zeta(2-s)\notag\\
 &=\frac{\pi\cot(\pi(s-1)/2)}{s-1}\zeta(s-1).
\end{align}
The second expression follows from the Riemann functional equation.
The completed germ is
\begin{equation}\label{beta:J1completion}
 J_1(1+u)+\frac1{u^2}+\frac L u.
\end{equation}
Its $u^m$ coefficient, multiplied by $(-1)^m m!$, is
$\FP\int\gamma_m\ell$. In particular,
\begin{equation}\label{beta:firstmoment}
 \FP\int_0^1\gamma_0(x)\ell(x)\dd x
 =\zeta''(0)+\frac{\zeta(2)}2.
\end{equation}
This also cross-checks the log-Gamma--cotangent specialization in
\cite{RepoExact}; neither \eqref{beta:J1} nor that specialization is
claimed to be new.

Write the \emph{inclusive} harmonic zeta function as
\begin{equation}\label{beta:Hzeta}
 \zeta_H(v)=\sum_{n\geq1}\frac{H_n}{n^v}
 =\frac1{(v-1)^2}+\frac\gamma{v-1}
  +\sum_{j\geq0}\frac{(-1)^j\gamma_H(j)}{j!}(v-1)^j.
\end{equation}
Here $\gamma_H(0)=(\gamma^2+\zeta(2))/2$; this normalization and
its Laurent theory are established in \cite{Kargin,Can}.

\begin{theorem}[The sine-square bridge]\label{beta:secondmoment}
One has
\begin{equation}\label{beta:J2}
 J_2(s)=B(s)\left\{L\zeta(2-s)+\zeta_H(2-s)
                       -\frac12\zeta(3-s)\right\}.
\end{equation}
Its completion at $s=1+u$ is
\begin{equation}\label{beta:J2completion}
 J_2(1+u)+\frac2{u^3}-\frac{L^2+\zeta(2)/2}{u}.
\end{equation}
In particular,
\begin{align}\label{beta:P2}
 \FP\int_0^1\psi(x)\ell(x)^2\dd x
 ={}&2\gamma_H(1)+\zeta'(2)+2L\gamma_1
    +\frac{\gamma^3}3+\gamma^2L\notag\\
   &-\frac{2L^3}3+\frac{2\zeta(3)}3.
\end{align}
Combining this with the already proved cubic harmonic bridge yields
the ordinary-constant identity
\begin{align}\label{beta:cancellation}
 \boxed{\begin{aligned}
 \FP\int_0^1\bigl\{\psi(x)^3-3\psi(x)\ell(x)^2\bigr\}\dd x
 ={}&3\zeta(2)+3\zeta'(2)-6(\gamma+L)\gamma_1\\
   &-\gamma^3-3\gamma^2L+2L^3-2\zeta(3).
 \end{aligned}}
\end{align}
\end{theorem}

\begin{proof}
Equation \eqref{beta:J2} is the second coefficient of
\eqref{beta:firstjets} in \eqref{beta:belltransform}.
The polar terms of \eqref{beta:J2completion} also follow directly
by taking the second $\lambda$ coefficient of
$C_0(\lambda)/u+1/(2\lambda-u)$ in
Theorem~\ref{beta:completion}. To make the constant calculation
explicit, set $a=\gamma+L$. Then
\[
 B(1+u)=-\frac2u
 \exp\left(au-\frac{\zeta(2)}4u^2+\frac{\zeta(3)}3u^3+O(u^4)\right),
\]
and the braces in \eqref{beta:J2} equal
\[
 \frac1{u^2}-\frac a u+L\gamma+\frac{\gamma^2}2
  +u\left(L\gamma_1+\gamma_H(1)+\frac{\zeta'(2)}2\right)+O(u^2).
\]
Multiplication gives the negative of \eqref{beta:P2}, because
$\gamma_0=-\psi$.

The incoming cubic report \cite{RepoCubic} proves, in precisely this
coordinate convention,
\begin{equation}\label{beta:inheritedcube}
 \FP\int_0^1\psi(x)^3\dd x
 =3\zeta(2)+6\gamma_H(1)+6\zeta'(2)-6\gamma\gamma_1.
\end{equation}
Subtracting three times \eqref{beta:P2} cancels $\gamma_H(1)$
identically and gives \eqref{beta:cancellation}. This use of the cubic
bridge is explicit: we prove the new cancellation, not a new proof
claim for the inherited bridge.
\end{proof}

The individual quantities are, for orientation,
\begin{align*}
 \FP\int_0^1\psi(x)^3\dd x
   &=1.9662627343915343406643753384116314\ldots,\\
 \FP\int_0^1\psi(x)\ell(x)^2\dd x
   &=-3.0645068706816349567044046185265270\ldots.
\end{align*}
Their combination in \eqref{beta:cancellation} is
$11.1597833464364392107775891939912123\ldots$.
These decimals come from independent endpoint-subtracted quadratures
and harmonic-tail evaluation in the supplied check script. They are
diagnostics; the coefficient identities and the inherited theorem are
the proof. No reduction of $\gamma_H(1)$ itself to ordinary constants
is inferred from its disappearance in this particular combination.
